\chapter{Standards Adopted}

\section{Design Standards}
\begin{itemize}
\item \textbf{IEEE Std 830 / ISO/IEC/IEEE 29148} -- structure of the requirements specification in Chapter~3.
\item \textbf{REST conventions} -- resource-style paths under \texttt{/api}, use of HTTP methods and JSON bodies, and meaningful status codes (200, 400, 413, 502).
\item \textbf{JSON (RFC 8259)} -- the only data-interchange format between browser, server and model.
\item \textbf{Responsive and accessible web design} -- semantic HTML, ARIA labels and live regions, native dialogs.
\end{itemize}

\section{Coding Standards}
\begin{itemize}
\item \textbf{PEP 8} naming and layout for Python modules, with type hints and docstrings in the sandbox and research modules.
\item Single-responsibility modules inside the \texttt{recall} package, with relative imports.
\item Configuration through environment variables (twelve-factor style); no secrets in source.
\item Defensive error handling at system boundaries: network calls have time-outs and failures degrade gracefully.
\end{itemize}

\section{Testing Standards}
\begin{itemize}
\item Black-box functional testing of every endpoint, including negative cases (invalid input, unsupported files, time-outs) as recorded in Table~\ref{tab:tests}.
\item Each test case states the input, the expected result and the observed result.
\item Recommended next step: convert the executed cases into an automated \texttt{pytest} suite using Flask's test client, with the LLM call mocked.
\end{itemize}
